\chapter[Covering laws and parapuzzles]{Source digests: covering laws, unicritical local connectivity, and parapuzzles}
\label{ch:digests-covering-parapuzzle}

This chapter continues Chapter~\ref{ch:digests-apriori-limbs}.  It digests three sources that
supply analytic and combinatorial inputs to the a priori bounds programme: the Kahn--Lyubich
quasi-additivity law and covering lemma, which are the analytic tools behind several of the
a priori bounds papers; the Kahn--Lyubich theorem on local connectivity of Julia sets of
unicritical polynomials; and Lyubich's Asterisque volume on parapuzzles and SBR measures, which
contains a local connectivity theorem for the Mandelbrot set under a hypothesis on families of
copies called \emph{fineness}.  The conventions are those of the earlier digest chapters.
Statements are the sources' own.  The Kahn--Lyubich laws are printed without numbers, so we
name them and do not invent numbers.

\section{Quasi-additivity law and covering lemma}
\label{sec:kl-qa}

\paragraph{Identity.}
J.~Kahn and M.~Lyubich, \emph{The quasi-additivity law in conformal geometry},
arXiv:math/0505191 (2005), Ann.\ of Math.\ 169 (2009), 561--593, DOI
10.4007/annals.2009.169.561 \cite{KahnLyubichQA}.  The title and authors are confirmed by the main
file \texttt{refs/arxiv-math-0505191/main.tex}, lines 17--22.  The proofs are in
\texttt{covering.tex}.

\paragraph{Definitions.}
\begin{itemize}
\item \emph{Extremal length and width.}  For a family $\Gamma$ of curves on a Riemann surface $U$ and
a conformal metric $\mu$ on $U$, put $\mu(\Gamma)=\inf_{\gamma\in\Gamma}\mu(\gamma)$ and
$\mathcal L_\mu(\Gamma)=\mu(\Gamma)^2/\mu^2(U)$.  The extremal length is
$\mathcal L(\Gamma)=\sup_\mu\mathcal L_\mu(\Gamma)$, and the extremal width is
$\mathcal W(\Gamma)=\mathcal L(\Gamma)^{-1}$ (\texttt{covering.tex}, lines 1658--1688).
\item \emph{Moduli of a compact set.}  For a compact $K\Subset\operatorname{int}S$ in a surface $S$,
$\mathcal L(S,K)=\modulus(S,K)$ is the extremal length of the curves joining $\partial S$ to $K$, and
$\mathcal W(S,K)=\mathcal L(S,K)^{-1}$ (\texttt{covering.tex}, lines 32--38 and 1825--1840).
\item \emph{Archipelagos and islands.}  An open archipelago is an open subset of $S$ whose closure is a
finite-type Riemann surface with smooth boundary; its components are its islands (\texttt{covering.tex},
lines 40--43).  For disjoint archipelagos $A_j$ the topological complexity is
$\operatorname{Top}_S\{A_j\}=-\chi(S)+\sum_j\#\operatorname{Comp}\,\partial A_j$ (lines 44--48).
\item \emph{The three moduli.}  For disjoint archipelagos $A_1,\dots,A_N$ in $S$ put
\[
X=\mathcal W\Bigl(S,\bigcup_jA_j\Bigr),\qquad
Y=\sum_j\mathcal W(S,A_j),\qquad
Z=\sum_j\mathcal W\Bigl(S\setminus\bigcup_{k\neq j}A_k,\,A_j\Bigr).
\]
The source notes the relations $X\le Y\le Z$ and $Y\le NX$ (\texttt{covering.tex}, lines 50--59).
\item \emph{Separation.}  The archipelagos are $\xi$-separated if $Z\le\xi Y$ (\texttt{covering.tex},
lines 91--94).
\end{itemize}

\paragraph{Results.}  The source prints the laws below as unnumbered statements.

\begin{theorem}[Quasi-additivity law]
\label{thm:kl-qa-law}
There exists $K$, depending only on the topological complexity of the family of archipelagos, such that
\[
Y\ \ge\ K\qquad\Longrightarrow\qquad Y^2\ \le\ 2\,XZ.
\]
\emph{Status:} proved in the paper; the proof occupies most of it (\texttt{covering.tex}, lines 60--70).
The source remarks that the constant $2$ may be replaced by any $C>4/3$, and that $C<32/27$ would not work
(lines 70--73).
\end{theorem}

\begin{theorem}[Quasi-additivity law with separation]
\label{thm:kl-qa-sep}
Assume that the archipelagos $A_j\Subset\operatorname{int}S$ are $\xi$-separated.  Then there is $K$,
depending only on $\xi$ and the topological complexity, such that
\[
Y\ \ge\ K\qquad\Longrightarrow\qquad Y\ \le\ 2\,\xi\,X.
\]
\emph{Status:} an immediate corollary of Theorem~\ref{thm:kl-qa-law} in the paper
(\texttt{covering.tex}, lines 91--110).
\end{theorem}

\begin{theorem}[Covering lemma]
\label{thm:kl-covering}
Fix $\eta\in(0,1]$.  Let $U\supset\Lambda'\supset\Lambda$ and $V\supset B'\supset B$ be nests of Jordan disks,
and let $f\colon(U,\Lambda',\Lambda)\to(V,B',B)$ be a branched covering with $D=\deg(U\to V)$ and
$d=\deg(\Lambda'\to B')$.  Suppose the collar assumption
\[
\modulus(B'\setminus B)\ >\ \eta\,\modulus(U\setminus\Lambda)
\]
holds.  Then there is $\epsilon=\epsilon(\eta,D)>0$ such that
\[
0<\modulus(U\setminus\Lambda)<\epsilon
\qquad\Longrightarrow\qquad
\modulus(V\setminus B)\ <\ 2\,\eta^{-1}d^2\,\modulus(U\setminus\Lambda).
\]
\emph{Status:} proved in the paper (\texttt{covering.tex}, lines 120--151).  A general estimate gives
$\modulus(V\setminus B)\le D\,\modulus(U\setminus\Lambda)$ without the collar assumption, which is why the
collar assumption is what makes the two moduli comparable.
\end{theorem}

\begin{proposition}[Quasi-invariance and holomorphic collars]
\label{prop:kl-quasi-invariance}
The quasi-additivity comparison of Theorem~\ref{thm:kl-qa-law} survives a branched covering of degree $D$,
in a form involving the degrees $d_j$ of the restrictions to the islands; the constant depends on the
topological complexity of the target archipelagos and on $D$.  The paper also proves the covering lemma with a
\emph{holomorphic} collar that excludes the critical values, with the same bound as in
Theorem~\ref{thm:kl-covering}.  \emph{Status:} proved in the paper; the precise formulations are in
\texttt{covering.tex}, lines 1415--1433 (quasi-invariance) and 1632--1655 (holomorphic collar).
\end{proposition}

\paragraph{Proof strategy (from the source).}
Leaves of a harmonic foliation are encoded by routes and associated rectangles (\texttt{covering.tex}, lines
175--226).  Buffers and a small-overlapping principle control the overlaps (lines 943--1029).  Truncated
rectangles are bounded by $Z$, and the initial ones by $X$ (lines 1030--1100).  The series law and an arithmetic
inequality give $(1-\delta)^2Y^2\le\frac43XZ$ (lines 1206--1264).  For the quasi-invariance law, one passes to a
Galois cover: a degree-$N$ branched covering factors through a Galois branched covering of degree at most $N!$
(lines 1351--1359), and the archipelagos are lifted and compared (lines 1351--1510).

\paragraph{External inputs.}
Ahlfors, \emph{Conformal Invariants} (1973), for extremal length and the Dirichlet-integral facts used
(\texttt{covering.tex}, lines 1658--1663 and 1841--1848).  The coset-action argument for the Galois cover is
elementary and is not cited externally (lines 1360--1380).  Later dynamical applications are to Kahn (2006),
Kahn--Lyubich (2005/2006), Kozlovski--van Strien (2006), and Qiu--Yin (2006) (\texttt{bib.tex}, lines 6--25).

\paragraph{Use in the MLC literature.}
The introduction says that the covering lemma is used to extend Yoccoz's theorem to higher-degree unicritical
maps, and to prove a priori bounds, hence local connectivity of $\M$, for some infinitely renormalizable quadratic
maps (\texttt{covering.tex}, lines 8--16).  Specifically, the basic covering lemma is used in the unicritical paper of
Section~\ref{sec:kl-unicritical}; the holomorphic-collar version in Kahn--Lyubich II; and the quasi-invariance law
with critical values inside the marked archipelagos in Kahn I (\texttt{covering.tex}, lines 1650--1655).
\emph{Scope:} the paper proves analytic laws only.  It does not claim MLC, and it says that the dynamical
consequences are addressed in forthcoming notes (\texttt{main.tex}, lines 23--42; \texttt{covering.tex},
lines 8--16).

\section{Local connectivity of Julia sets for unicritical polynomials}
\label{sec:kl-unicritical}

\paragraph{Identity.}
J.~Kahn and M.~Lyubich, \emph{Local connectivity of Julia sets for unicritical polynomials},
arXiv:math/0505194 (2005), Ann.\ of Math.\ 170 (2009), 413--426, DOI 10.4007/annals.2009.170.413
\cite{KahnLyubichUnicritical}.  The title and authors are confirmed by the main file
\texttt{refs/arxiv-math-0505194/main.tex}, lines 19--20.

\paragraph{Definitions.}
\begin{itemize}
\item \emph{Unicritical polynomials.}  The maps are $f\colon z\mapsto z^d+c$ with $d\ge2$ (\texttt{intro.tex},
lines 6--12).  No other normalization is used.
\item \emph{GPL maps.}  A generalized polynomial-like map is a holomorphic map $g\colon\bigcup_iW_i\to V$, where $V$
is a topological disk, the $W_i\Subset V$ are disjoint topological disks, every restriction $g|_{W_i}$ is a
branched covering, and all but finitely many of these restrictions have degree one (\texttt{kingdom.tex},
lines 4--9).  A GPL map is \emph{unicritical} if it has one critical point, of arbitrary local degree (\texttt{intro.tex},
lines 72--73).
\item \emph{Puzzle pieces.}  The puzzle pieces of depth $n$ are the components of $g^{-n}(V)$; those containing the
critical point $0$ are critical puzzle pieces (\texttt{kingdom.tex}, lines 39--41).
\item \emph{Principal nest.}  A sequence of critical puzzle pieces $V^0\supset V^1\supset\cdots\ni0$, chosen inductively
so that the first-return maps $f^{n_i}\colon V^i\to V^{i-1}$ are unicritical branched coverings (\texttt{intro.tex},
lines 43--46).
\item \emph{Modified principal nest.}  A nest $E^0\supset E^1\supset\cdots\ni0$ of associated domains, built from kingdom
maps with enlargements and buffers (\texttt{intro.tex}, lines 15--19; \texttt{kingdom.tex}, lines 426--438).  A kingdom
map has nested disks $W\supset U\supset A\ni0$ together with subjects and men, and it is ramified only at $0$
(\texttt{kingdom.tex}, lines 67--84).
\item \emph{At most finitely renormalizable.}  After finitely many satellite or primitive renormalizations one reaches a
non-renormalizable quadratic-like map (\texttt{kingdom.tex}, lines 287--295).
\end{itemize}

\paragraph{Results.}

\begin{theorem}[Local connectivity for unicritical polynomials]
\label{thm:kl-unicritical-A}
The Julia set $J(f)$ of an at most finitely renormalizable unicritical polynomial $f\colon z\mapsto z^d+c$ with all
periodic points repelling is locally connected.  \emph{Status:} proved in the paper, for every unicritical degree $d$
(\texttt{intro.tex}, lines 1--13).
\end{theorem}

\begin{theorem}[A priori bounds in the modified principal nest]
\label{thm:kl-unicritical-B}
The modified principal moduli stay away from zero:
\[
\modulus(E^{i-1}\setminus E^i)\ \ge\ \mu\ >\ 0.
\]
\emph{Status:} proved in the paper (\texttt{intro.tex}, lines 17--24), via Lemma~\ref{lem:kl-growth} below.  The source
then says that these bounds force the puzzle pieces $E^i$ to shrink to the critical point, which yields
Theorem~\ref{thm:kl-unicritical-A} by a standard argument (\texttt{intro.tex}, lines 25--27).
\end{theorem}

\begin{lemma}[Growth lemma]
\label{lem:kl-growth}
There exist $n=n(d)\in\N$ and $\epsilon=\epsilon(d,n)>0$ such that the following holds.  If on some level $q\ge n$
\[
\modulus(V^q\setminus W^q)<\epsilon,
\]
then on some previous level $p<q$
\[
\modulus(V^p\setminus W^p)\ <\ \tfrac12\,\modulus(V^q\setminus W^q).
\]
\emph{Status:} proved in the paper (\texttt{bounds.tex}, lines 4--14).  Its proof uses the Quasi-Additivity Law applied to
many disjoint domains; this is the decisive step (\texttt{bounds.tex}, lines 45--68).
\end{lemma}

\begin{lemma}[Time inequalities]
\label{lem:kl-times}
The travel times $t_n$ and the auxiliary times $T_n$, $s_n$ of the source satisfy $t_n\ge T_{n-1}$, $T_n\ge s_n$ and
$s_n\ge2s_{n-1}$ (\texttt{kingdom.tex}, lines 490--507).  Consequently $t_n=\operatorname{Time}(W^n,V^n)\ge2^{n-1}$
(Corollary ``growth'', \texttt{kingdom.tex}, lines 535--539).  \emph{Status:} proved in the paper.
\end{lemma}

\paragraph{Proof strategy (from the source).}
The argument reduces a finitely renormalizable polynomial to a non-renormalizable member of its renormalization
dynasty (\texttt{kingdom.tex}, lines 287--295).  It builds the modified principal nest from kingdom maps (lines
103--152).  Exponential growth of return times and the time inequality give bounded-degree pullbacks
(\texttt{intro.tex}, lines 52--59).  The Covering Lemma of Section~\ref{sec:kl-qa} transfers small-modulus estimates
between levels (\texttt{intro.tex}, lines 63--77), and the Quasi-Additivity Law gives the growth lemma.

\paragraph{External inputs.}
Yoccoz's quadratic local connectivity theorem, cited through Hubbard and Milnor (\texttt{intro.tex}, lines 6--8).
Lyubich, \emph{Dynamics of quadratic polynomials I--II}, for the principal-nest and puzzle machinery (\texttt{intro.tex},
lines 43--49).  The Quasi-Additivity Law and the Covering Lemma of Kahn--Lyubich (Section~\ref{sec:kl-qa}), which the paper
states as inputs and does not reprove (\texttt{covering.tex}, lines 7--32).  Branner--Hubbard (puzzle machinery, 1992),
Douady--Hubbard (polynomial-like maps, 1985), and Levin--van Strien (local connectivity of real polynomials, 1998)
(\texttt{intro.tex}, lines 38--40 and 79--81; \texttt{bib.tex}).

\paragraph{MLC scope.}
Theorem~\ref{thm:kl-unicritical-A} is a theorem about Julia sets, not about $\M$.  For $d=2$ and at most finitely
renormalizable parameters it reproduces Yoccoz's theorem, which is the known part of MLC (Chapter~\ref{ch:digests-foundations}).
The paper contains no statement that $\M$ is locally connected, and it does not claim the MLC conjecture.

\section{Parapuzzles and SBR measures}
\label{sec:lyubich-ast261}

\paragraph{Identity.}
M.~Lyubich, \emph{Dynamics of quadratic polynomials, III: parapuzzle and SBR measures}, Asterisque 261 (2000),
173--200 \cite{LyubichIII}.  The article is in the Asterisque archive at NUMDAM, \texttt{numdam.org/item?id=AST\_2000\_\_261\_\_173\_0};
the statements below were checked against that PDF, and page references are to its journal pages.  The PDF is not
stored in \texttt{refs/}; its URL is the canonical location.

\paragraph{Definitions.}
\begin{itemize}
\item \emph{SBR measure.}  A measure $\mu$ is SBR (Sinai--Bowen--Ruelle) if
\[
\frac1n\sum_{k=0}^{n-1}f^k_*\delta_x\ \longrightarrow\ \mu
\]
for $x$ in a set of positive Lebesgue measure (the source's equation (1.1)).  For a real quadratic map $P_c$, $c\in[-2,1/4]$,
the SBR measure is taken on the invariant interval $I_c$.
\item \emph{Quadratic-like and DH families.}  A quadratic-like family over a parameter disk $D$ is a fibrewise
quadratic-like family with a holomorphic motion respecting the boundary dynamics.  A DH family has a single tube $V_0$
(pp.~175--176).  For such a family $f$ the
Mandelbrot set of the family is $\M(f)=\{\lambda\in D:\ J(f_\lambda)\text{ is connected}\}$ (p.~176).
\item \emph{Bounded shape.}  A copy $M'$ of the Mandelbrot set has $(K,\epsilon)$-\emph{bounded shape} if its straightening
extends to a $K$-quasiconformal map on the $(\epsilon\operatorname{diam}M')$-neighbourhood of $M'$ (p.~174).
\item \emph{Fine family.}  A family $\mathcal C$ of Mandelbrot copies is \emph{fine} if all copies of each specified
combinatorial type in $\mathcal C$ have uniformly bounded shape (p.~198).  The examples given are the family of
maximal copies with big height, and the family of real maximal copies with big essential period (p.~199).
\item \emph{Parapuzzle pieces and principal nest.}  The first tiling is made of the sets
$A_1(\lambda)=\{\lambda\in O:\ f^p_\lambda(0)\in L^*_\lambda\}$, where $O$ is a Misiurewicz wake and $L^*_\lambda$ is the tile
containing the critical value (p.~186).  The principal parapuzzle nest $A^0\supset A^1(\lambda)\supset A^2(\lambda)\supset\cdots$
is finite if and only if $\lambda$ is renormalizable, and its principal parameter moduli are
$\modulus(A^l(\lambda)\setminus A^{l+1}(\lambda))$ (p.~191).
\end{itemize}

\paragraph{Results.}

\begin{theorem}[Parameter moduli grow linearly]
\label{thm:ast-A}
Let $(f,h)$ be a proper DH quadratic-like family with winding number $1$ over $D$, and let $O\subset D$ be a Misiurewicz
wake.  Then for every $\lambda\in\M(f)\cap O$ and every $l$,
\[
\modulus\bigl(A^l(\lambda)\setminus A^{l+1}(\lambda)\bigr)\ >\ Bl,
\]
where $B>0$ depends only on $O$ and on $\modulus(f)$ (p.~191).  The source states a second
inequality of the same type, $\modulus(A^l(\lambda)\setminus H^l(\lambda))>Bl$, with a set $H^l(\lambda)$ whose symbol the OCR
renders imperfectly.  \emph{Status:} proved in the paper; it is the main geometric result of its Section~4.
\end{theorem}

\begin{theorem}[Theorem C: bounded shape of Mandelbrot copies]
\label{thm:ast-C}
For every Misiurewicz wake $O$, the maximal Mandelbrot copies contained in $O$ have bounded shape.  In particular all maximal
real Mandelbrot copies, except the doubling one, have bounded shape.  Moreover a real copy $M'$ has $(K,\epsilon)$-bounded shape
with $K\to1$ and $\epsilon\to\infty$ as the essential period $p_e(M')\to\infty$ (p.~175).
\emph{Status:} proved in the paper.
\end{theorem}

\begin{theorem}[Theorem B, joint with Martens and Nowicki]
\label{thm:ast-B}
Lebesgue almost every real quadratic map $P_c(z)=z^2+c$ which is non-hyperbolic and at most finitely renormalizable has a
finite absolutely continuous invariant measure (p.~174).  \emph{Status:} proved in the paper
from Theorem~\ref{thm:ast-A} (through Theorem~5.1 below), the Martens--Nowicki summability criterion (Section~\ref{sec:martens-nowicki}), and Lyubich's exponential
decay of scaling factors in the quasi-quadratic case \cite{LyubichQuasiQuadratic}.
\end{theorem}

\begin{corollary}[SBR measures for almost every real parameter]
\label{cor:ast-sbr}
For almost all $c\in[-2,1/4]$ the quadratic polynomial $P_c$ has a unique SBR measure on its invariant interval $I_c$
(p.~174).  \emph{Status:} derived in the paper from Theorem~\ref{thm:ast-B}.  The
the uniqueness of the SBR measure is the Martens--Nowicki Theorem~\ref{thm:mn-A}, stated jointly with Lyubich (Section~\ref{sec:martens-nowicki}).
\end{corollary}

\begin{theorem}[Theorem 5.1: real non-renormalizable parameters]
\label{thm:ast-5-1}
Let $\mathcal M_{\mathbb R}\subset[-2,-3/4)$ be the non-renormalizable real parameters, and let $\mathcal M_C\subset\mathcal M_{\mathbb R}$
be those whose principal nest contains only finitely many nontrivial central cascades.  Then $\mathcal M_{\mathbb R}$ has positive
Lebesgue measure, and $\mathcal M_C$ has full measure in $\mathcal M_{\mathbb R}$ (p.~196).
\emph{Status:} proved in the paper; the positive-measure part is known from Benedicks--Carleson and Jakobson, and the full-measure
part is new in the source.
\end{theorem}

\begin{proposition}[Local connectivity on fine families]
\label{prop:ast-fine}
Let $\mathcal C$ be a fine family of Mandelbrot copies, and let $\mathcal E_{\mathcal C}$ be the set of infinitely renormalizable
parameters whose combinatorics $(M_0,M_1,\dots)$ satisfy $M_n\in\mathcal C$.  Then $\M$ is locally connected at every point of
$\mathcal E_{\mathcal C}$, and $\mathcal E_{\mathcal C}$ has zero Lebesgue measure (pp.~198--199).
\emph{Status:} proved in the paper, conditional on the fineness of $\mathcal C$.  The topological step of the proof (nested copies with diameters tending to zero meet in one point) is formalized as \leanref{Mlc/Literature/NestedCopies.lean}{exists_unique_mem_iInter_of_antitone}; bounded shape and fineness are not formalized.
\end{proposition}

\paragraph{Proof strategy (from the source).}
Holomorphic motions and transverse quasiconformality transfer the dynamical puzzle geometry to parameter space.  The initial
parapuzzle tiling is built in Misiurewicz wakes from critical-value itineraries (pp.~185--186).  Each tile is renormalized
inductively to give the principal parapuzzle nest (p.~191).  Parameter annuli are compared with dynamical annuli by uniformly
quasiconformal pseudo-conjugacies, and linear modulus growth gives Theorem~\ref{thm:ast-A}.  For the real line, the central
cascades are estimated, and Borel--Cantelli together with the Martens--Nowicki criterion gives Theorem~\ref{thm:ast-B}.

\paragraph{External inputs.}
Lyubich, \emph{Dynamics of quadratic polynomials I--II} (Acta Math.\ 178 (1997)), quoted for a modulus estimate in
Section~4 and for the essential period (\cite{LyubichDynamics}).  Lyubich, \emph{Combinatorics, geometry and attractors of quasi-quadratic maps}
(Ann.\ of Math.\ 140 (1994)) \cite{LyubichQuasiQuadratic}.  Martens and Nowicki, \emph{Invariant measures for Lebesgue typical
quadratic maps}, in the same Asterisque volume \cite{MartensNowicki}.  Douady--Hubbard, \emph{On the dynamics of polynomial-like maps}
(1985) \cite{DHPolynomialLike}.  Douady's ICM report \emph{Chirurgie sur les applications holomorphes} (1986) \cite{DouadyICM}, for the
winding-number fact used in Section~3.  Benedicks--Carleson (Ann.\ of Math.\ 122 (1985)) \cite{BenedicksCarleson} and Jakobson
(positive measure of real parameters; cited in the source as Jacobson [J]).  Holomorphic motions and quasiconformal extension:
Slodkowski, Man\'e--Sad--Sullivan, Sullivan--Thurston, and Bers--Royden, as cited in the source.

\paragraph{MLC relevance.}
The paper proves local connectivity of $\M$ only at points of $\mathcal E_{\mathcal C}$ for a fine family $\mathcal C$ of
Mandelbrot copies (Proposition~\ref{prop:ast-fine}).  It does not prove MLC.  The author conjectures that the family of real maximal
copies is fine, which would give bounded shape for all real copies.  Theorem~\ref{thm:ast-C} gives bounded shape for the maximal real copies other than the doubling one, but fineness is only conjectured.  He remarks that
fineness of the full family of complex maximal copies would imply MLC, and says that he is not sure whether this holds (p.~199).  Julia-set local connectivity enters only as cited background.

\section{The Martens--Nowicki criterion}
\label{sec:martens-nowicki}

\paragraph{Identity.}
M.~Martens and T.~Nowicki, \emph{Invariant measures for typical quadratic maps}, arXiv:math/9606220, Asterisque 261
(2000), 239--252 \cite{MartensNowicki}.  The title and authors are confirmed by the metadata and by
\texttt{refs/arxiv-math-9606220/main.tex}, lines 15--31; the journal pages are those listed by NUMDAM, since the
source files do not give them.

\paragraph{Definitions.}
\begin{itemize}
\item \emph{$S$-unimodal maps.}  A map $f\colon[-1,1]\to[-1,1]$ is $S$-unimodal with critical exponent $\alpha$ if it is
$C^3$, $f(\pm1)=-1$, $Df>0$ on $x<0$, $Df<0$ on $x>0$, $Df(0)=0$, $Sf<0$, and near $0$ it is $x\mapsto|x|^\alpha$ up to a
differentiable change of coordinates, with $\alpha>1$ (\texttt{main.tex}, lines 838--849).  The quadratic family is
$q_t(x)=-2tx^2+2t-1$ (lines 67--72).
\item \emph{Nice points and first returns.}  $\mathcal N=\{x:\ f^i(x)\notin(-|x|,|x|)\text{ for all }i\ge0\}$ (lines 253--258).
For $x\in\mathcal N$ put $U_x=(-|x|,|x|)$; the points of $U_x$ that return to $U_x$ form $D_x$, and the first-return map is
$R_x\colon D_x\to U_x$ (lines 260--264).  The central component $V_x\subset D_x$ contains $0$, and $\psi(x)=\partial V_x\cap(0,1)$
(lines 286--305).
\item \emph{Scaling factors.}  Choose $u_1\in\mathcal N$, put $u_n=\psi(u_{n-1})$ and $U_n=(-u_n,u_n)$; then $U_{n+1}$ is the
central component of $D_n$ (lines 306--317).  The scaling factor of level $n$ is $\sigma_n=u_{n+1}/u_n$ (lines 319--328).
\end{itemize}

\paragraph{Results.}

\begin{theorem}[Martens--Nowicki, Theorem B]
\label{thm:mn-B}
If an $S$-unimodal map with critical exponent $\alpha>1$ has summable scaling factors,
\[
\sum_{n\ge1}\sigma_n^{1/\alpha}\ <\ \infty,
\]
then it has an absolutely continuous invariant probability measure (\texttt{main.tex}, lines 191--199).  \emph{Status:}
proved in the paper, using the Nowicki--van Strien summability criterion \cite{NowickiVanStrien} (lines 633--645).
\end{theorem}

\begin{corollary}[Martens--Nowicki, Corollary]
\label{cor:mn-finite}
A quadratic map with only finitely many central returns has an absolutely continuous invariant probability measure
(\texttt{main.tex}, lines 201--204).  \emph{Status:} derived from Theorem~\ref{thm:mn-B} together with a renormalization
input that the paper cites.
\end{corollary}

\begin{theorem}[Martens--Nowicki, Theorem A, joint with Lyubich]
\label{thm:mn-A}
A typical quadratic map has a unique SBR measure.  More precisely, for parameters in $\mathcal P$ the support of the SBR
measure is the periodic attractor; for parameters in $\mathcal R$ it is a Cantor set; for parameters in $\mathcal M$ it is
absolutely continuous and supported on the orbit of a periodic interval; and $\mathcal P\cup\mathcal R\cup\mathcal M$ has
full Lebesgue measure (\texttt{main.tex}, lines 121--145).  \emph{Status:} stated in the paper as joint with Lyubich.  The
paper derives it from Theorem~\ref{thm:mn-B} and from Lyubich's result that for almost every parameter in $\mathcal I$ the
map has only finitely many central returns (lines 215--223).
\end{theorem}

\paragraph{Proof strategy (from the source).}
\begin{itemize}
\item Proposition~3.1: small scaling factors give derivative growth along the critical orbit (\texttt{main.tex}, lines 386--401).
\item The Nowicki--van Strien summability criterion turns this growth into an absolutely continuous invariant measure
(lines 633--645).
\item A Teleman-type decomposition of the critical orbit splits the derivative sum; the outside pieces are estimated
exponentially by Proposition~4.1 (cited from Guckenheimer and Ma\~n\'e) and the return pieces by Proposition~3.1 (lines 661--724).
\item Lemma~4.2 shows that distinct decompositions cannot have the same data, so the products decay geometrically and the sum
converges (lines 725--763).
\end{itemize}

\paragraph{External inputs.}
Nowicki--van Strien, \emph{Absolutely continuous measures under a summability condition}, Invent.\ Math.\ 93 (1988)
\cite{NowickiVanStrien}; Martens, \emph{Interval dynamics} (thesis, 1990), for distortion (\texttt{main.tex}, lines 179--181);
Lyubich, \emph{Combinatorics, geometry and attractors of quasi-quadratic maps} \cite{LyubichQuasiQuadratic}, for the exponential
decay of scaling factors (lines 181--187); Lyubich--Milnor, \emph{The Fibonacci unimodal map}, J.\ Amer.\ Math.\ Soc.\ 6 (1993),
425--457 (lines 181--187); Guckenheimer and Ma\~n\'e, for Proposition~4.1; and Lyubich's parapuzzle preprint, cited as [L4] in the
paper's bibliography (\texttt{Dynamics of Quadratic Polynomials, III: Parapuzzles}, Stony Brook preprint).

\paragraph{Relation to Lyubich's Theorem B.}
Lyubich's Astérisque paper uses the criterion in the following form: together with the exponential decay of scaling factors
in the quasi-quadratic case, the criterion gives an absolutely continuous invariant measure once the principal nest is eventually
free from central cascades (Section~\ref{sec:lyubich-ast261}).  Theorem~\ref{thm:ast-5-1} supplies this condition for a set of full
measure among the non-renormalizable real parameters.  The bibliography entry [L4] of the Martens--Nowicki paper has the title of
the parapuzzle paper of Part~III; we identify it with the Astérisque paper by this title series, which is our identification and
not a statement of the source.

\paragraph{Scope.}
The paper does not claim that every non-hyperbolic quadratic map has an absolutely continuous invariant measure; it notes that
for parameters in $\mathcal I$ large enough there are maps without one (\texttt{main.tex}, lines 213--216).  The uniqueness statement
of Theorem~\ref{thm:mn-A} concerns typical parameters only.

\section{Summary of this chapter}
\label{sec:digests-iv-summary}

\begin{table}[ht]
\centering\small
\begin{tabularx}{\textwidth}{@{}>{\raggedright\arraybackslash}p{0.2\textwidth}>{\raggedright\arraybackslash}X>{\raggedright\arraybackslash}p{0.27\textwidth}@{}}
\textbf{Source} & \textbf{Result (label in this chapter)} & \textbf{Status in the source}\\
\hline
Kahn--Lyubich, quasi-additivity &
Quasi-additivity law; law with separation; covering lemma (Theorems~\ref{thm:kl-qa-law}--\ref{thm:kl-covering}) &
proved (the laws are unnumbered in the source)\\
Kahn--Lyubich, unicritical &
Local connectivity of Julia sets of at most finitely renormalizable unicritical polynomials (Theorem~\ref{thm:kl-unicritical-A}) &
proved for every degree $d$; no statement about $\M$\\
Lyubich, Asterisque 261 &
Linear growth of parameter moduli (Theorem~\ref{thm:ast-A}); bounded shape of copies (Theorem~\ref{thm:ast-C}) &
proved\\
Lyubich, Asterisque 261 &
$\M$ locally connected on $\mathcal E_{\mathcal C}$ for a fine family (Proposition~\ref{prop:ast-fine}) &
proved, conditional on fineness\\
\end{tabularx}
\caption{Results digested in this chapter and their status in the sources.}
\label{tab:digests-iv-summary}
\end{table}

\paragraph{Logical role.}
The covering laws of Section~\ref{sec:kl-qa} are the analytic tools of the Kahn--Lyubich unicritical theorem (Section~\ref{sec:kl-unicritical}) and of the a priori bounds series of Chapter~\ref{ch:digests-apriori-limbs}.  In the parapuzzle paper, the parameter analogue of the same moduli growth gives local connectivity of $\M$ on fine families; the fineness hypothesis is the point where that paper's argument meets the open problem.

\paragraph{Not covered here.}
Martens and Nowicki's criterion for absolutely continuous invariant measures is cited, not digested.  Lyubich's 1994 Annals paper \cite{LyubichQuasiQuadratic} and the Acta papers \cite{LyubichDynamics} are cited, not digested in this chapter.  The Lean development formalizes only the topological step of Proposition~\ref{prop:ast-fine} (nested copies with shrinking diameters meet in one point, \leanref{Mlc/Literature/NestedCopies.lean}{subsingleton_iInter_of_diam_tendsto_zero}); the statements of this chapter otherwise have no Lean counterpart.
